\documentclass{article} % For LaTeX2e
\usepackage{iclr2025_conference,times}

\usepackage{hyperref}
\usepackage{url}


\title{Learning Developer Preferences from In-the-Wild Coding Agent Logs}

\author{Haowen Wang \\
Stanford University \\
\texttt{haowenw@stanford.edu} \\
\And
Christy Yang \\
Stanford University \\
\texttt{yangyx@stanford.edu} \\
\And
Zhen Jia \\
Stanford University \\
\texttt{zhenjia1@stanford.edu}
}

\iclrfinalcopy

\begin{document}


\maketitle

\section{Problem Specification and Research Questions}

Coding agents start each session without a memory of the developer, so they repeat behavior that the developer already corrected in earlier sessions.
We study whether a forecaster can learn a developer's preferences from past sessions in one repository and use them to predict the developer's reaction to a new agent turn.
We use SWE-chat \citep{baumann2026swechat} and keep only repositories with one or two contributors, so that each repository stands for one developer.
Each input is one target turn: the conversation so far and the coding agent's replies and tool calls since the last developer message.
The forecaster also gets read-only access to the first $k=4$ sessions of the same repository, which all end before the target session starts.
The output is the probability that the developer's next typed message is pushback.
We define pushback as a reaction to an avoidable error of the coding agent: a failure report, a rejection of its work or a tool call, or a correction of a misunderstood or ignored instruction.
New requirements, context that the agent could not know, and answers to the agent's questions are not pushback.
This definition is narrower than SWE-chat's labels, whose \textit{correction} class also covers these cases.
SWE-chat labels 53\% of its 111{,}034 target turns as pushback, but our labeler marks only 51 of 233 random English targets (22\%) as pushback.
On a pilot of 27 hand-labeled turns, SWE-chat's binary label matched our hand labels on only 15.
We ask three questions.
\textbf{RQ1:} Does repository history improve the forecast over a model that sees only the target turn, and over the repository's past pushback rate?
\textbf{RQ2:} Is history more useful as a memory of verified preference rules than as raw past episodes found by BM25 retrieval?
A rule becomes active only when exact quotes from two independent past sessions support it.
\textbf{RQ3:} Is the gain specific to the developer, so that memory from the same repository beats memory from a random other repository?
The forecaster succeeds if it beats the two RQ1 baselines and the other-repository memory on AUROC, with 95\% bootstrap confidence intervals over repositories.
A forecaster that meets this criterion can later serve as a personalized reward model for training a coding agent with GRPO.

\section{Environment and Data}

\paragraph{Environment.}
We reorganize the raw SWE-chat interaction logs \citep{baumann2026swechat} into an offline forecasting environment. For each target turn, the forecaster observes the conversation context and the coding agent's replies and tool calls since the latest developer message. The memory agent can inspect eligible historical interactions through read-only search and episode-inspection tools. The target's subsequent developer feedback is withheld and used only for labeling and evaluation. The system predicts pushback without modifying a live software repository.

\paragraph{Historical and evaluation data.}
We construct examples from SWE-chat. Repository identity serves as an imperfect proxy for developer identity. Each historical episode from the first $k=4$ sessions contains the developer request, agent activity, subsequent developer feedback, and relabeled outcome.

Our development evaluation uses the k4-100 sample, which contains 100 targets: 20 pushback, 73 non-pushback, and 7 marked as not feedback. Excluding the last category leaves 93 forecastable targets from 50 repositories. We additionally relabel the 1{,}089 historical turns associated with this sample. Labels are generated by a separate GPT-6.1-sol relabeling pipeline using our operational definition of pushback. They are model-generated annotations rather than human-verified ground truth.

\paragraph{Input boundaries.}
We separate forecasting inputs from evaluation annotations. The target's subsequent developer message and outcome label are withheld from the forecaster. Preferences and explanations produced during relabeling are also excluded from forecasting inputs. Historical developer feedback remains available because it precedes the target and provides evidence from which the agent must infer preferences. These chronological constraints prevent access to future feedback; they do not by themselves eliminate every possible shortcut in recorded agent activity.

\paragraph{Evaluation.}
We evaluate discrimination using AUROC and AUPRC, and probability quality using Brier score. We also report accuracy and F1 at a threshold of 0.5. Because pushback is a minority class, we include an always-non-pushback accuracy baseline and do not judge success using accuracy alone. We report valid-output coverage and compare methods on shared valid targets. We use repository-level bootstrap intervals for selected paired comparisons to account for dependence among targets from the same repository. Current samples are development data; final claims require evaluation with a frozen configuration on an untouched test sample.

\section{Methods}

\paragraph{Forecasting task.}
Let $x_t$ denote the current conversation context and coding-agent activity, and let $H_r$ denote the eligible historical interactions for repository $r$. The forecaster estimates
\begin{equation}
p_t = \Pr(y_t = 1 \mid x_t, H_r),
\end{equation}
where $y_t=1$ indicates that the next developer message expresses pushback under our definition. The approaches below differ in how they represent and access $H_r$.

\paragraph{Baselines.}
B0 predicts the pushback rate of the training targets for every target. B1 predicts the pushback rate in the repository's first $k$ sessions and falls back to B0 when these sessions have no labeled turns. B0 and B1 are implemented but currently use SWE-chat labels rather than our relabeled targets. B2 is a direct language-model predictor receiving only $x_t$. B3 augments B2 with the five historical episodes ranked highest by BM25. Retrieved episodes include the historical request, agent activity, developer feedback, and outcome label. This tests whether retrieving raw examples is sufficient without extracting explicit preferences.

\paragraph{ReAct forecaster.}
We implement the agent using the DSPy framework \citep{khattab2023dspy}, with language-model modules for preference extraction, merging, verification, and applicability selection. A ReAct controller \citep{yao2023react} inspects the current interaction and, when available, historical evidence before returning a pushback probability and supporting citations. Its tools support historical search and inspection of individual episodes. Citation checks validate references to available evidence, although they do not establish that the evidence justifies the prediction. We also evaluate the agent framework without historical memory to distinguish the effect of history from that of the controller.

\paragraph{Preference memory.}
Memory is built once per repository from its first $k$ sessions and shared by all targets of that repository. The memory pipeline processes historical interactions in batches, extracts candidate preference rules, merges related candidates, and verifies their supporting evidence. Each rule stores its text, supporting episode references and quotes, confidence, and verification status. A candidate rule is promoted to active status only when exact supporting quotes are verified across at least two independent historical sessions. The original session-memory variant also exposes candidate summaries to the forecaster.

Our scoped variant addresses the risk of generalizing a one-time instruction into a persistent preference. It distinguishes task-specific instructions from stated preferences and recurring requirements, excluding task-specific instructions and items judged non-reusable from reusable memory. Each retained rule includes an applicability condition, which is preserved during merging. Before forecasting, a separate selection step compares active rules with the current developer request. Only rules judged applicable, with a supporting quote from the current request, are supplied to the forecaster; missing or invalid applicability evidence causes a rule to be withheld. This conservative gate may reduce interference from unrelated preferences but may also suppress useful rules. We therefore examine extraction, activation, and selection separately when interpreting forecasting performance.

\paragraph{Planned personalization ablation.}
To test whether any benefit is specific to the developer or repository, we will compare same-repository memory with memory drawn from another repository. Together with the no-history and BM25 baselines, this ablation will help distinguish personalized preferences from generic historical examples.


\section{Summary of Progress}

From SWE-chat's 229{,}909 user prompts, we built 111{,}034 target turns, of which 81{,}978 come from the 645 repositories with one or two contributors.
We split the 519 of these repositories that have at least two sessions into train, holdout, and test sets (70/15/15), so that no user appears in two sets.
Our relabeling pipeline applies the pushback definition of Section~1 and makes sure that every evidence quote occurs in the episode.
It matched our hand labels on 24 of 27 pilot turns, but we tuned the prompt on these turns.
Besides k4-100 (Section~2), we labeled a balanced $k=4$ set, k4-30 (15 pushback, 15 no pushback).
In DSPy, we implemented the session memory, its scoped variant, the ReAct forecaster with citation checks, baselines B0--B3 and 131 offline tests.

\paragraph{Preliminary results.}
Of the 93 forecastable k4-100 targets, 91 have a valid output from every method: B3 failed the citation check on 1 target, and scoped memory construction failed on 1 target.
On these 91 targets, neither memory variant beats simple retrieval (Table~\ref{tab:k4-100}).
Scoped memory lowers the ReAct forecaster's Brier score by 0.055 (95\% repository bootstrap CI $[0.017, 0.094]$), but session memory lowers it by only 0.023 ($[-0.014, 0.060]$).
No scoped rule reached the forecaster, so the scoped gain likely comes from the forecaster's direct access to raw history episodes rather than from preference rules.
B3 has the best score on every metric.
Its Brier score is 0.056 lower than session memory's ($[0.022, 0.095]$), and its difference from scoped memory is not significant ($[-0.003, 0.055]$).
No method meets the AUROC criterion of Section~1: all AUROC values are near 0.5, and the three ReAct variants assign a lower mean probability to pushback targets than to non-pushback targets.
All methods also overestimate pushback (mean prediction 0.33--0.41, true rate 0.22).
On the 29 targets of the balanced k4-30 set where every method produced a valid output, the session-memory agent (AUROC 0.69) and B3 (0.70) beat B2 (0.66), but all confidence intervals include 0.

\begin{table}[t]
\caption{k4-100 results on the 91 targets where every method produced a valid output (22\% pushback). Forecasts use GPT-6-luna. F1 and accuracy use a 0.5 threshold. Always predicting no pushback gives an accuracy of 0.780.}
\label{tab:k4-100}
\begin{center}
\begin{tabular}{lccccc}
\multicolumn{1}{c}{\bf METHOD} &\multicolumn{1}{c}{\bf AUROC} &\multicolumn{1}{c}{\bf AUPRC} &\multicolumn{1}{c}{\bf BRIER $\downarrow$} &\multicolumn{1}{c}{\bf F1} &\multicolumn{1}{c}{\bf ACC.}
\\ \hline \\
B2: LLM, no history     &0.457 &0.239 &0.257 &0.311 &0.659 \\
B3: BM25 + LLM          &\bf 0.530 &\bf 0.293 &\bf 0.230 &\bf 0.326 &\bf 0.681 \\
ReAct, no history       &0.430 &0.206 &0.309 &0.259 &0.560 \\
ReAct + session memory (Ours)  &0.460 &0.215 &0.286 &0.245 &0.593 \\
ReAct + scoped memory (Ours)   &0.455 &0.257 &0.255 &0.222 &0.615 \\
\end{tabular}
\end{center}
\end{table}


\paragraph{Next steps.}
We will hand-label 100 targets and report Cohen's $\kappa$.
We will also run B1 on the relabeled targets and build the other-repository memory ablation for RQ3.
Scoped memory yields only 23 rules across 49 repositories, and only 2 become active; we will audit rule extraction and cross-session merging to find out why.
We will calibrate with Platt scaling on holdout.
Then we will freeze the configuration and evaluate it once on a test-split sample with at least 60 pushback targets.

\paragraph{Obstacles.}
The history can contain little signal beyond the current turn.
If so, a tie with B1 or B3 is itself a finding.
With 20 pushback targets, most confidence intervals include 0.
Larger runs are expensive: the paired ReAct run with and without memory used 949 model calls for 93 targets (about 10 per target), while B2 and B3 use one call each.
One LLM produced all the labels, and we inspected the holdout set many times during development.
Memory construction still fails for one k4-100 history set.
An AUROC near 0.5 is too weak for a reward model.
If the criterion fails, we will replace GRPO with best-of-$N$ reranking or drop that stage.

\subsubsection*{AI Disclosure}
We used AI coding assistants for two parts of this project.
\textbf{Demo page.} Claude Code helped write the demo page: the HTML template.
\textbf{Code.} Claude Code helped write parts of the pipeline code, including [the DSPy memory and ReAct forecaster modules, the baselines, the evaluation scripts, and the offline tests].
We reviewed, ran, and tested all code ourselves and take responsibility for it.

\subsubsection*{Github Code Link}
Code: \url{https://github.com/whw-alex/CS-329Z-Project-public}

\bibliography{iclr2025_conference}
\bibliographystyle{iclr2025_conference}

\end{document}
